% Part: incompleteness
% Chapter: representability-in-q
% Section: comp-representable

\documentclass[../../../include/open-logic-section]{subfiles}

\begin{document}

\olfileid{inc}{req}{crq}
% PNB-CMP-060: converse تے stronger-theory مثالاں لئی ضروری موافقت دا فرض صاف۔
\olsection{قابلِ حساب فنکشن $\Th{Q}$ وچ قابلِ نمائندگی نیں}

\begin{thm}
ہر قابلِ حساب فنکشن~$\Th{Q}$ وچ قابلِ نمائندگی اے۔
\end{thm}

\begin{proof}
گل نوں ٹھیک طرح متعین کرن لئی چرچ--ٹیورنگ دا دعویٰ ورتیے۔
اسیں کسے فنکشن نوں اُدوں تے صرف اُدوں قابلِ حساب آکھاں گے
جدوں اوہ عمومی بازگشتی ہووے۔ عمومی بازگشتی فنکشناں دی تعریف صفر
فنکشن~$\Zero$،جانشین فنکشن~$\Succ$ تے پروجیکشن
فنکشن~$\Proj{n}{i}$ توں ترکیب، بنیادی بازگشت تے باقاعدہ کمترین قدر دی تلاش نال
کیتی جا سکدی اے۔\olref[pri]{lem:prim-rec} نال،ہر فنکشن~$h$
جیہڑا $f$ تے~$g$ توں بنیادی بازگشت نال تعریف کیتا جا سکدا اے،اوہ
ترکیب تے باقاعدہ کمترین قدر دی تلاش نال $f$، $g$ تے $\Zero$،
$\Succ$، $\Proj{n}{i}$، $\Add$، $\Mult$، $\Char{=}$ توں وی تعریف کیتا جا سکدا اے۔ایس لئی کوئی فنکشن اُدوں تے صرف اُدوں عمومی بازگشتی اے جدوں اوہ $\Zero$،
$\Succ$، $\Proj{n}{i}$، $\Add$، $\Mult$، $\Char{=}$ توں ترکیب
تے باقاعدہ کمترین قدر دی تلاش نال تعریف کیتا جا سکدا اے۔

اسیں ایہ وی وکھا چکے آں کہ متعلقہ بنیادی فنکشن
~$\Th{Q}$ وچ قابلِ نمائندگی نیں
(\cref{inc:req:bre:prop:rep-zero,inc:req:bre:prop:rep-succ,inc:req:bre:prop:rep-proj,inc:req:bre:prop:rep-id,inc:req:bre:prop:rep-add,inc:req:bre:prop:rep-mult})،
تے قابلِ نمائندگی فنکشناں توں
ترکیب یا باقاعدہ کمترین قدر دی تلاش نال تعریف کیتا گیا ہر فنکشن
(\olref[cmp]{prop:rep-composition}،
\olref[min]{prop:rep-minimization}) وی قابلِ نمائندگی اے۔ ایس لئی ہر
عمومی بازگشتی فنکشن~$\Th{Q}$ وچ قابلِ نمائندگی اے۔
\end{proof}

\begin{explain}
اسیں وکھایا اے کہ قابلِ حساب فنکشناں دا سیٹ اوہناں سارے
فنکشناں دا سیٹ اے جیہڑے $\Th{Q}$ وچ قابلِ نمائندگی نیں۔ ایہ ثبوت ہور وی عام
اے۔ نمائندگی دی تعریف توں سوکھا ای دِسدا اے کہ ہر اوہ نظریہ
جیہڑا $\Th{Q}$ نوں ودھاؤندا اے یا جیہدے وچ $\Th{Q}$ دی تعبیر
ہو سکدی اے، قابلِ حساب فنکشناں دی نمائندگی کر سکدا اے۔ اُلٹی
سمت وچ، ایہو جیہے کسے بے تضاد نظریے دے !!{derivation} نظام وچ
جے !!{derivation} ہون دی پرکھ قابلِ حساب اے، تاں ہر قابلِ
نمائندگی فنکشن قابلِ حساب اے۔ بے تضادی دا ایہ فرض رکھ کے،
قابلِ حساب فنکشناں دا سیٹ پیانو حساب یا زرمیلو--فرینکل سیٹ
نظریے وچ قابلِ نمائندگی فنکشناں دا سیٹ وی اے۔ جیویں گوڈل نے
دَسیا، ایہ گل ذرا حیران کرن والی اے۔ اسیں ویکھاں گے کہ کیہ
ثابت ہو سکدا اے، ایہ چُنے نظریے اُتے بہت انحصار کردا اے۔
عام طور اُتے مسلّمے جِنّے طاقتور ہون، اوِنّا ودھ ثابت ہو سکدا
اے۔ پر مسلّماتی نظرییاں دے اک وڈے دائرے وچ قابلِ نمائندگی
فنکشن عین قابلِ حساب فنکشن نیں۔ جے اوہ نظریے بے تضاد تے مسلّماتی
طور متعین ہون، تاں ہور طاقتور نظریے ہور فنکشناں دی نمائندگی
نہیں کردے۔
\end{explain}

\end{document}
